\section{Summary of HARL integration into MAPush repository}
\label{app:harl_integration}

This appendix documents the integration of HAPPO (from the HARL library)
into MAPush, supporting all three task levels (mid-level, velocity, and
upper-level planner) with homogeneous and heterogeneous robot teams.

% -----------------------------------------------------------
\subsection{Architecture}
% -----------------------------------------------------------

The integration wraps MAPush as a HARL environment while keeping
MAPush-specific training and testing code in a separate
\texttt{harl\_mapush/} directory, minimising changes to the core HARL
library.

\begin{dirtree}{%
.1 HARL/.
.2 harl/.
.3 envs/mapush/.
.4 mapush\_env.py \DTcomment{env wrapper, $\sim$1000 lines}.
.4 mapush\_logger.py \DTcomment{TensorBoard logging}.
.3 utils/envs\_tools.py \DTcomment{mapush registration}.
.2 harl\_mapush/.
.3 train.py \DTcomment{40+ CLI flags, 3 tasks}.
.3 test.py \DTcomment{calculator, viewer, batch, video}.
.3 runners/mapush\_happo\_runner.py \DTcomment{step-based checkpoints}.
}

Only two core HARL files were modified: \texttt{harl/envs/\_\_init\_\_.py}
(logger registration) and \texttt{harl/utils/envs\_tools.py} (environment
factory).

% -----------------------------------------------------------
\subsection{Environment wrapper}
% -----------------------------------------------------------

\texttt{mapush\_env.py} adapts Isaac~Gym's MAPush to HARL's interface.
It dispatches environment creation based on robot composition
(\texttt{make\_mqe\_env} for homogeneous Go1, \texttt{make\_hetero\_env}
for all other pairings) and dispatches \texttt{step()}/\texttt{reset()}
based on task type.

\paragraph{Task dispatch.}
Table~\ref{tab:task-dispatch} summarises the three supported tasks.

\begin{table}[h]
  \centering
  \caption{Task dispatch in the HARL environment wrapper.}
  \label{tab:task-dispatch}
  \begin{tabular}{lcccc}
    \toprule
    \textbf{Task} & \textbf{Agents} & \textbf{Obs dim} & \textbf{Act dim} & \textbf{Dispatch} \\
    \midrule
    \texttt{go1push\_mid}   & 2 & 8  & 3 & \texttt{\_step\_mid()} \\
    \texttt{go1push\_vel}   & 2 & 15 & 3 & \texttt{\_step\_vel()} \\
    \texttt{go1push\_upper} & 1 & 26 & 2 & \texttt{\_step\_upper()} \\
    \bottomrule
  \end{tabular}
\end{table}

For the upper-level task, HARL sees a single agent ($n=1$). The upper
wrapper internally handles both robots via frozen mid-level actors.

\paragraph{Centralised critic modes.}
Five global state representations are available for the centralised critic
(Table~\ref{tab:critic-modes}), selectable via CLI flags. The state type
is always EP (Environment Provided). For the velocity task, a 15-dim
box-centred state is used by default (18-dim legacy variant via
\texttt{--legacy\_vel\_obs}). For the upper task, the 26-dim observation
already contains all global information, so it is reused as the critic state.

\begin{table}[h]
  \centering
  \caption{Centralised critic state modes (mid-level task).}
  \label{tab:critic-modes}
  \begin{tabular}{llp{7cm}}
    \toprule
    \textbf{Mode} & \textbf{Dims} & \textbf{Description} \\
    \midrule
    Absolute (default)         & 11 & box(3) + target(2) + agent0(3) + agent1(3) \\
    Box-centred                & 9  & target\_rel(2) + agents\_rel(3 each) + box\_yaw(1) \\
    Goal-centred               & 9  & box\_rel(3) + agents\_rel(3 each) \\
    Relative                   & 9  & robots\_to\_box(3 each) + inter\_dist(1) + goal\_to\_box(2) \\
    Concat-agent               & 16 & agent0\_obs(8) + agent1\_obs(8) \\
    \bottomrule
  \end{tabular}
\end{table}

% -----------------------------------------------------------
\subsection{Training and testing scripts}
% -----------------------------------------------------------

The training script (\texttt{train.py}) exposes 40+ CLI flags grouped into:
core (\texttt{--algo}, \texttt{--task}, \texttt{--seed}),
training (\texttt{--n\_rollout\_threads}, \texttt{--num\_env\_steps}),
heterogeneous agents (\texttt{--agent0}, \texttt{--agent1}),
box physics (\texttt{--box\_mass}, \texttt{--box\_mass\_range}),
reward shaping (\texttt{--contact\_force\_gating}, etc.),
critic mode selection,
velocity-task parameters, and
upper-task parameters (\texttt{--mid\_level\_checkpoint}, \texttt{--mid\_level\_format}).

The custom runner (\texttt{MAPushHAPPORunner}) extends HARL's
\texttt{OnPolicyHARunner} with step-based checkpoints (every 10M steps),
disabled HARL evaluation (Isaac~Gym single-instance limitation), and
automatic saving of \texttt{config.json}, \texttt{command.txt}, and
\texttt{run\_config.yaml} at training start. Each checkpoint folder
(e.g.\ \texttt{10M/}) contains \texttt{actor\_agent\{i\}.pt},
\texttt{critic\_agent.pt}, and \texttt{value\_normalizer.pt}.

The testing script (\texttt{test.py}) provides three modes:
\textbf{calculator} (multi-environment statistics),
\textbf{viewer} (sequential visualisation), and
\textbf{batch} (\texttt{--all\_checkpoints} tests 10M\ldots200M and
outputs a summary table).
Video recording is available via \texttt{--record\_video}.

% -----------------------------------------------------------
\subsection{Upper-level task integration}
% -----------------------------------------------------------

The upper-level planner trains a single HARL agent that outputs sub-goal
waypoints $(x, y)$, which are executed by frozen mid-level actors:

\begin{equation*}
  \underbrace{\text{Upper (HAPPO, 1 agent)}}_{\text{26D obs} \to \text{2D sub-goal}}
  \;\to\;
  \underbrace{\text{Mid-level (frozen, 2 actors)}}_{\text{8D obs} \to [v_x, v_y, v_\text{yaw}]}
  \;\to\;
  \underbrace{\text{Locomotion}}_{\text{joint torques}}
\end{equation*}

\paragraph{Integration work.}
Six files were modified to support this pipeline:

\begin{enumerate}
  \item \textbf{Upper wrapper} (\texttt{go1\_push\_upper\_wrapper.py}): added
        \texttt{\_prepare\_happo\_command\_policy()}, which loads per-agent
        \texttt{StochasticPolicy} networks from HARL checkpoints with matching
        architecture (hidden sizes, activation, initialisation). A new
        \texttt{\_happo\_mid\_level\_act()} method runs inference on both
        frozen actors under \texttt{torch.no\_grad()}, returning scaled
        velocity commands. The original OpenRL loading path is retained,
        selected at runtime via \texttt{cfg.control.mid\_level\_format}.

  \item \textbf{Upper config} (\texttt{go1\_push\_upper\_config.py}): added
        \texttt{mid\_level\_format = "openrl"} as backward-compatible default.

  \item \textbf{Config pipeline} (\texttt{mqe/envs/utils.py}): extended
        \texttt{custom\_cfg()} to accept and propagate
        \texttt{mid\_level\_checkpoint} and \texttt{mid\_level\_format}.

  \item \textbf{Training script} (\texttt{train.py}): added
        \texttt{go1push\_upper} to task choices, and
        \texttt{--mid\_level\_checkpoint} / \texttt{--mid\_level\_format}
        CLI flags.

  \item \textbf{Environment wrapper} (\texttt{mapush\_env.py}): added upper-task
        detection (\texttt{is\_upper\_task}), space overrides
        (obs=26D, act=2D, $n_\text{agents}=1$), and dedicated
        \texttt{\_step\_upper()} / \texttt{\_reset\_upper()} methods that
        convert between HARL tensor shapes and the upper wrapper's interface.
        NaN sanitisation was added at both step and reset to handle physics
        corruption from degenerate quaternions.

  \item \textbf{Testing script} (\texttt{test.py}): extended all three modes
        for the upper task---single-agent checkpoint detection, mid-level
        config passthrough, upper-specific metrics (distance-to-target,
        trajectory, reach-target, obstacle rewards), and correct action
        tensor shapes for single-agent inference.
\end{enumerate}

\paragraph{Usage.}

\begin{lstlisting}[style=bashstyle, language=bash]
# Train upper-level planner with frozen HAPPO mid-level
python HARL/harl_mapush/train.py \
    --task go1push_upper \
    --mid_level_checkpoint results/mapush/.../checkpoints/190M/ \
    --mid_level_format happo \
    --n_rollout_threads 10 --num_env_steps 200000000

# Test upper-level planner
python HARL/harl_mapush/test.py \
    --checkpoint results/mapush/.../checkpoints/10M \
    --task go1push_upper \
    --mid_level_checkpoint results/mapush/.../checkpoints/190M/ \
    --mode calculator --num_envs 50 --num_episodes 100
\end{lstlisting}

% -----------------------------------------------------------
\subsection{Heterogeneous robot support}
% -----------------------------------------------------------

Three robot types are supported via \texttt{--agent0}/\texttt{--agent1}
flags. All expose a unified 3-action interface $[v_x, v_y, v_\text{yaw}]$
so HARL sees homogeneous action spaces regardless of robot heterogeneity.

\begin{table}[h]
  \centering
  \caption{Supported robot types.}
  \label{tab:robots}
  \begin{tabular}{lllll}
    \toprule
    \textbf{Robot} & \textbf{Type} & \textbf{DOFs} & \textbf{Locomotion} & \textbf{Actuator} \\
    \midrule
    go1      & Quadruped & 12 & Walk-These-Ways (70D + 2100 history) & Feedforward \\
    anymal\_c & Quadruped & 12 & legged\_gym MLP (48D) & LSTM \\
    cassie   & Biped     & 12 & legged\_gym MLP (48D) & PD control \\
    \bottomrule
  \end{tabular}
\end{table}

The routing logic directs \texttt{go1+go1} to \texttt{make\_mqe\_env()}
(native path) and all other combinations to \texttt{make\_hetero\_env()},
backed by \texttt{HeteroRobot} (1444 lines: multi-URDF loading, per-agent
locomotion policies, per-agent actuator networks, action routing).

% -----------------------------------------------------------
\subsection{Contact-force gating}
% -----------------------------------------------------------

To combat freeloading in heterogeneous teams, rewards can be gated by the
balance of contact forces between agents.  Isaac~Gym's
\texttt{net\_contact\_force\_tensor} provides per-body horizontal forces;
these are summed over non-foot bodies and normalised by agent mass:

\begin{equation}
  c_i = \frac{F_{\text{horiz},i}}{m_i}, \qquad
  b = \frac{\min_i\, c_i}{\max_i\, c_i}, \qquad
  g = \alpha + (1-\alpha)\, b
  \label{eq:cf-gate}
\end{equation}

The gate $g$ is applied to \texttt{push\_reward} and
\texttt{distance\_to\_target\_reward} when box speed exceeds 0.01\,m/s.
Flags: \texttt{--contact\_force\_gating True --contact\_force\_gating\_alpha 0.3}.

% -----------------------------------------------------------
\subsection{File summary}
% -----------------------------------------------------------

\begin{table}[h]
  \centering
  \caption{Files created and modified for the integration.}
  \label{tab:files}
  \begin{tabular}{p{7cm}lp{4.5cm}}
    \toprule
    \textbf{File} & \textbf{Lines} & \textbf{Purpose} \\
    \midrule
    \multicolumn{3}{l}{\textit{Created (HARL-side)}} \\
    \texttt{harl/envs/mapush/mapush\_env.py}               & $\sim$1000 & Environment wrapper (3 tasks) \\
    \texttt{harl/envs/mapush/mapush\_logger.py}             & $\sim$145  & TensorBoard logging \\
    \texttt{harl\_mapush/train.py}                          & $\sim$290  & Training entry point \\
    \texttt{harl\_mapush/test.py}                           & $\sim$830  & Testing (all modes, 3 tasks) \\
    \texttt{harl\_mapush/runners/mapush\_happo\_runner.py}  & $\sim$240  & Custom runner \\
    \midrule
    \multicolumn{3}{l}{\textit{Modified (HARL core -- minimal)}} \\
    \texttt{harl/envs/\_\_init\_\_.py}                       & +1 & Logger registration \\
    \texttt{harl/utils/envs\_tools.py}                      & +12 & Environment factory \\
    \midrule
    \multicolumn{3}{l}{\textit{Modified (MAPush-side)}} \\
    \texttt{mqe/envs/wrappers/go1\_push\_upper\_wrapper.py} & +125 & HAPPO mid-level actor loading \\
    \texttt{mqe/envs/configs/go1\_push\_upper\_config.py}   & +1   & \texttt{mid\_level\_format} default \\
    \texttt{mqe/envs/utils.py}                              & +8   & Mid-level config passthrough \\
    \bottomrule
  \end{tabular}
\end{table}
